\documentclass[12pt, a4paper]{article}
\usepackage[utf8]{inputenc}
\usepackage[T1]{fontenc}
\usepackage[french]{babel}
\usepackage[margin=2.5cm]{geometry}
\usepackage{amsmath}
\usepackage{xcolor}
\usepackage{tcolorbox}
\tcbuselibrary{theorems, skins, breakable}
\usepackage{booktabs}
\usepackage{fancyhdr}
\usepackage{hyperref}
\usepackage{enumitem}
\usepackage{tikz}
\usetikzlibrary{arrows.meta, positioning, shapes.geometric, decorations.pathreplacing}
\usepackage{array}
\usepackage{multirow}
\usepackage{longtable}
\usepackage{colortbl}

% Couleurs
\definecolor{suva_blue}{RGB}{30, 100, 180}
\definecolor{done_green}{RGB}{34, 139, 34}
\definecolor{done_bg}{RGB}{230, 255, 230}
\definecolor{inprog_orange}{RGB}{200, 100, 0}
\definecolor{inprog_bg}{RGB}{255, 245, 220}
\definecolor{todo_red}{RGB}{180, 30, 30}
\definecolor{todo_bg}{RGB}{255, 235, 235}
\definecolor{critical_bg}{RGB}{255, 200, 200}
\definecolor{phase_blue}{RGB}{200, 220, 255}
\definecolor{warn_bg}{RGB}{255, 248, 220}

% Boîtes
\newtcolorbox{done}[2][]{colback=done_bg,colframe=done_green,
  fonttitle=\bfseries\color{done_green},title={\checkmark\ #2},#1,breakable}
\newtcolorbox{inprog}[2][]{colback=inprog_bg,colframe=inprog_orange,
  fonttitle=\bfseries\color{inprog_orange},title={$\triangleright$\ #2},#1,breakable}
\newtcolorbox{todo}[2][]{colback=todo_bg,colframe=todo_red,
  fonttitle=\bfseries\color{todo_red},title={$\square$\ #2},#1,breakable}
\newtcolorbox{critical}[2][]{colback=critical_bg,colframe=todo_red!80!black,
  fonttitle=\bfseries,title={!! #2 !!},#1,breakable}
\newtcolorbox{notebox}[1][]{colback=suva_blue!5,colframe=suva_blue!50,
  fonttitle=\bfseries,title={Note},#1}

% Commandes de statut
\newcommand{\statdone}{\textcolor{done_green}{\bfseries FAIT}}
\newcommand{\statprog}{\textcolor{inprog_orange}{\bfseries EN COURS}}
\newcommand{\stattodo}{\textcolor{todo_red}{\bfseries À FAIRE}}
\newcommand{\statna}{\textcolor{gray}{\bfseries N/A}}

\pagestyle{fancy}
\fancyhf{}
\fancyhead[L]{\color{suva_blue}\textbf{SuVa --- Roadmap}}
\fancyhead[R]{\color{gray}État d'avancement du projet}
\fancyfoot[L]{\footnotesize\color{gray}Mohamed Lamine MBENGUE}
\fancyfoot[C]{\thepage}
\fancyfoot[R]{\footnotesize\color{gray}portfolio-alamine.web.app}
\renewcommand{\headrulewidth}{0.4pt}
\renewcommand{\footrulewidth}{0.3pt}

% ══════════════════════════════════════════════════════════════════════
\begin{document}

\begin{center}
  {\color{suva_blue}\rule{\linewidth}{2pt}}\\[0.5cm]
  {\Huge\bfseries\color{suva_blue} Projet SuVa}\\[0.3cm]
  {\large\bfseries État d'avancement \& Feuille de route}\\[0.2cm]
  {\small\color{gray} Stablecoin quantique-résistant sur EVM --- Mars 2026}\\[0.3cm]
  {\color{suva_blue}\rule{\linewidth}{2pt}}
\end{center}

\vspace{0.3cm}
\begin{tcolorbox}[colback=suva_blue!5, colframe=suva_blue!40,
  left=6pt, right=6pt, top=4pt, bottom=4pt]
\small
\begin{tabular}{ll}
  \textbf{Auteur}      & Mohamed Lamine MBENGUE \\
  \textbf{Spécialités} & Sécurité · DevOps · Dev Full Stack Web \& Mobile \\
  \textbf{Contact}     & +221 78 178 44 36 · mbenguemohamedlamine2022@gmail.com \\
  \textbf{Portfolio}   & \href{https://portfolio-alamine.web.app}{\texttt{portfolio-alamine.web.app}} \\
\end{tabular}
\end{tcolorbox}
\vspace{0.3cm}

\tableofcontents
\newpage

%══════════════════════════════════════════════════════════════════════
\section{Lexique --- Définitions des termes clés}
%══════════════════════════════════════════════════════════════════════

\begin{tcolorbox}[colback=suva_blue!4, colframe=suva_blue!60,
  title=\textbf{Comment utiliser ce lexique},
  fonttitle=\bfseries, breakable]
Ce document utilise beaucoup de termes techniques propres à la blockchain et à la cryptographie.
Ce lexique les explique simplement, \textbf{sans prérequis}. Consulte-le à chaque fois qu'un
terme te semble flou.
\end{tcolorbox}

\vspace{0.4cm}

%──────────────────────────────────────────────────────────────────────
\subsection{Termes Blockchain de base}
%──────────────────────────────────────────────────────────────────────

\begin{tcolorbox}[colback=gray!4, colframe=gray!50, breakable]

\textbf{Blockchain} (chaîne de blocs)\\
\hspace*{1em} Une base de données \textit{partagée entre des milliers d'ordinateurs dans le monde},
où chaque donnée ajoutée est \textbf{permanente et publiquement vérifiable}. Personne ne peut
modifier ou supprimer ce qui a été écrit. C'est comme un registre public et inviolable.
\textit{Exemples : Ethereum, Bitcoin.}

\vspace{0.4cm}

\textbf{Ethereum}\\
\hspace*{1em} La blockchain la plus utilisée pour les applications (DeFi, NFT, etc.). Elle permet
d'exécuter du code (les \textit{smart contracts}) directement sur la chaîne.

\vspace{0.4cm}

\textbf{EVM} (Ethereum Virtual Machine --- Machine Virtuelle Ethereum)\\
\hspace*{1em} Le ``moteur'' qui exécute le code des smart contracts sur Ethereum. Tous les
ordinateurs participants font tourner ce même moteur, ce qui garantit que tout le monde obtient
le même résultat.

\vspace{0.4cm}

\textbf{Mainnet} (Réseau principal)\\
\hspace*{1em} La vraie blockchain Ethereum, avec de l'argent réel. Opposé au \textit{testnet}
(réseau de test) où l'on utilise de l'argent fictif pour tester le code sans risque.

\vspace{0.4cm}

\textbf{Testnet} (Réseau de test)\\
\hspace*{1em} Une copie de la blockchain Ethereum utilisée uniquement pour tester. Les tokens
n'ont aucune valeur réelle. \textit{Exemples : Sepolia, Holesky.}

\vspace{0.4cm}

\textbf{Layer 2 (L2)} (Couche 2)\\
\hspace*{1em} Une blockchain construite \textit{par-dessus} Ethereum pour aller plus vite et moins
cher. Les L2 regroupent des milliers de transactions et les soumettent en une seule fois à
Ethereum. \textit{Exemples : Polygon, Arbitrum, Optimism.}

\end{tcolorbox}

\vspace{0.4cm}

%──────────────────────────────────────────────────────────────────────
\subsection{Termes financiers et tokens}
%──────────────────────────────────────────────────────────────────────

\begin{tcolorbox}[colback=gray!4, colframe=gray!50, breakable]

\textbf{Stablecoin} (monnaie stable)\\
\hspace*{1em} Une cryptomonnaie dont la valeur est maintenue \textbf{fixe} (ex. : toujours 1 dollar).
Contrairement au Bitcoin dont le prix fluctue, un stablecoin est stable. C'est pratique pour
stocker ou transférer de la valeur sans risque de volatilité.\\
\hspace*{1em} \textit{Exemples : USDT (Tether), USDC (USD Coin) --- tous les deux valent ~1\$.}

\vspace{0.4cm}

\textbf{USDT / USDC}\\
\hspace*{1em} Les deux stablecoins les plus populaires, chacun valant environ 1 dollar américain.
\textbf{USDT} est émis par Tether, \textbf{USDC} par Circle. SuVa veut les ``wrapper'' (envelopper)
pour les rendre quantique-résistants.

\vspace{0.4cm}

\textbf{sUSDT / sUSDC} (le ``s'' = ``SuVa-secured'')\\
\hspace*{1em} Les versions sécurisées créées par SuVa. Quand tu déposes 100 USDT dans le vault SuVa,
tu reçois 100 sUSDT. Le sUSDT est garanti 1:1 par l'USDT d'origine, mais protégé par une
signature post-quantique.

\vspace{0.4cm}

\textbf{Token}\\
\hspace*{1em} Un actif numérique sur une blockchain. Un token peut représenter de l'argent,
un droit de vote, un accès à un service, etc. USDT et sUSDT sont des tokens.

\vspace{0.4cm}

\textbf{ERC-20}\\
\hspace*{1em} Le \textit{standard} (ensemble de règles) le plus courant pour créer des tokens sur
Ethereum. Tous les tokens ERC-20 fonctionnent de la même façon :
\texttt{transfer()}, \texttt{approve()}, \texttt{balanceOf()}... Cela permet aux wallets et
applications de les utiliser sans avoir à adapter leur code.

\end{tcolorbox}

\vspace{0.4cm}

%──────────────────────────────────────────────────────────────────────
\subsection{Termes de programmation blockchain}
%──────────────────────────────────────────────────────────────────────

\begin{tcolorbox}[colback=gray!4, colframe=gray!50, breakable]

\textbf{Smart Contract} (Contrat Intelligent)\\
\hspace*{1em} Un programme qui s'exécute \textbf{automatiquement} sur la blockchain, sans
intermédiaire. Une fois déployé, son code est public et immuable. Il exécute des règles
prédéfinies : ``\textit{si X envoie 100 USDT, alors créditer X de 100 sUSDT}''.

\vspace{0.4cm}

\textbf{Solidity}\\
\hspace*{1em} Le langage de programmation utilisé pour écrire des smart contracts sur Ethereum.
Ressemble un peu à JavaScript/C++. C'est ce qui est dans les fichiers \texttt{.sol} du projet.

\vspace{0.4cm}

\textbf{Yul}\\
\hspace*{1em} Un langage d'assembleur de bas niveau pour l'EVM, utilisé dans les parties critiques
des contrats pour \textbf{réduire le coût gas}. C'est l'équivalent de l'assembleur x86 mais pour
Ethereum. C'est difficile à lire mais très efficace.

\vspace{0.4cm}

\textbf{Vault} (Coffre-fort)\\
\hspace*{1em} Dans le contexte SuVa, le vault est le smart contract qui \textbf{garde les USDT}
déposés par les utilisateurs. Il les libère uniquement si la signature post-quantique est valide.
C'est le c\oe{}ur financier du protocole.

\vspace{0.4cm}

\textbf{Gas}\\
\hspace*{1em} L'\textbf{unité de mesure du coût} d'exécution d'une opération sur Ethereum.
Chaque instruction EVM a un coût en gas. Le gas se paie en ETH (l'ether, la monnaie native
d'Ethereum). Plus une opération est complexe, plus elle coûte de gas.\\
\hspace*{1em} \textit{Repères : une addition = 3 gas, un accès mémoire = 3 gas, une vérification
ECDSA = 3000 gas, une vérification Dilithium = 6 à 13 millions de gas (!)}.

\vspace{0.4cm}

\textbf{On-chain / Off-chain}\\
\hspace*{1em} \textbf{On-chain} = exécuté sur la blockchain, visible de tous, coûte du gas.\\
\hspace*{1em} \textbf{Off-chain} = exécuté localement sur ton ordinateur, gratuit, privé.\\
\hspace*{1em} Dans SuVa : la génération de clés et la signature se font off-chain (Python),
la vérification se fait on-chain (Solidity).

\vspace{0.4cm}

\textbf{Déploiement} (deploy)\\
\hspace*{1em} L'action de publier un smart contract sur la blockchain pour la première fois.
Une fois déployé, le contrat reçoit une adresse unique et son code est immuable.

\vspace{0.4cm}

\textbf{Nonce}\\
\hspace*{1em} Un numéro qui \textbf{augmente à chaque transaction} d'un compte. Il empêche
de réutiliser deux fois la même transaction (attaque par rejeu). Si le nonce est 5, la
prochaine transaction doit avoir le nonce 6.

\vspace{0.4cm}

\textbf{Replay Attack} (Attaque par rejeu)\\
\hspace*{1em} Une attaque où un malfaiteur copie une transaction légitime et la ``rejoue'' pour
voler des fonds. Le nonce l'empêche : une transaction avec un nonce déjà utilisé est rejetée.

\end{tcolorbox}

\vspace{0.4cm}

%──────────────────────────────────────────────────────────────────────
\subsection{Termes wallet et identité}
%──────────────────────────────────────────────────────────────────────

\begin{tcolorbox}[colback=gray!4, colframe=gray!50, breakable]

\textbf{Wallet} (Portefeuille numérique)\\
\hspace*{1em} Un logiciel qui gère tes clés cryptographiques et te permet d'interagir avec
la blockchain. Il ne ``stocke'' pas vraiment les tokens (ils sont sur la chaîne), il stocke
la \textbf{clé privée} qui prouve que tu es bien le propriétaire.
\textit{Exemples : MetaMask, Ledger, Trust Wallet.}

\vspace{0.4cm}

\textbf{ERC-4337} (Account Abstraction --- Abstraction de compte)\\
\hspace*{1em} Un standard Ethereum qui permet de remplacer le mécanisme de signature habituel
(ECDSA) par \textbf{n'importe quelle logique de validation personnalisée} --- notamment une
signature post-quantique. Au lieu d'un compte contrôlé par une clé ECDSA, l'utilisateur a
un \textit{smart contract wallet} qui valide les signatures avec Dilithium ou Falcon.

\vspace{0.4cm}

\textbf{ECDSA} (Elliptic Curve Digital Signature Algorithm)\\
\hspace*{1em} L'algorithme de signature utilisé par Ethereum aujourd'hui pour prouver qu'une
transaction est bien autorisée par le propriétaire du compte. \textbf{Vulnérable aux ordinateurs
quantiques} (algorithme de Shor). SuVa vise à le remplacer par Dilithium ou Falcon.

\vspace{0.4cm}

\textbf{Clé privée (sk) / Clé publique (pk)}\\
\hspace*{1em} La \textbf{clé privée} est secrète : elle te permet de signer des transactions.
Ne jamais la partager. La \textbf{clé publique} est partageable : elle permet à n'importe qui
de \textit{vérifier} ta signature sans connaître ta clé privée.

\vspace{0.4cm}

\textbf{Bundler}\\
\hspace*{1em} Dans ERC-4337, un bundler est un service qui \textbf{collecte les opérations des
utilisateurs} et les regroupe en une seule transaction Ethereum pour économiser du gas. C'est
l'équivalent d'un taxi collectif pour les transactions.

\vspace{0.4cm}

\textbf{Paymaster}\\
\hspace*{1em} Dans ERC-4337, un contrat optionnel qui \textbf{paie le gas à la place de
l'utilisateur}. Cela permet par exemple de payer les frais de transaction en sUSDT plutôt
qu'en ETH.

\vspace{0.4cm}

\textbf{Multi-sig} (Multi-signatures)\\
\hspace*{1em} Un wallet qui \textbf{requiert plusieurs signatures} pour valider une transaction.
Par exemple : ``3 personnes parmi 5 doivent signer''. Utile pour la gouvernance ou la
récupération de compte en cas de perte de clé.

\end{tcolorbox}

\vspace{0.4cm}

%──────────────────────────────────────────────────────────────────────
\subsection{Termes de sécurité et cryptographie}
%──────────────────────────────────────────────────────────────────────

\begin{tcolorbox}[colback=gray!4, colframe=gray!50, breakable]

\textbf{Signature numérique}\\
\hspace*{1em} L'équivalent numérique d'une signature manuscrite. Elle prouve que le détenteur de
la clé privée a bien approuvé un message ou une transaction. Elle est \textbf{vérifiable par
n'importe qui} avec la clé publique correspondante.

\vspace{0.4cm}

\textbf{Post-quantique (PQ)}\\
\hspace*{1em} Qualifie des algorithmes cryptographiques \textbf{résistants aux ordinateurs
quantiques}. Un ordinateur quantique suffisamment puissant pourrait casser ECDSA (la signature
d'Ethereum). Les algorithmes PQ (Dilithium, Falcon, SPHINCS+) sont basés sur des problèmes
mathématiques que même les ordinateurs quantiques ne savent pas résoudre efficacement.

\vspace{0.4cm}

\textbf{Ordinateur quantique}\\
\hspace*{1em} Un ordinateur utilisant les lois de la physique quantique pour effectuer certains
calculs \textbf{exponentiellement plus vite} qu'un ordinateur classique. Il ne cassera pas tout
(les bases de données, le chiffrement symétrique restent sûrs), mais il peut casser ECDSA et RSA.
Les experts estiment qu'un tel ordinateur pourrait exister dans 10 à 20 ans.

\vspace{0.4cm}

\textbf{Hash} (Empreinte numérique / Hachage)\\
\hspace*{1em} Une fonction qui transforme une donnée de taille quelconque en une empreinte de
taille fixe. C'est à sens unique : on ne peut pas retrouver la donnée d'origine à partir du hash.
\textit{Exemple : SHA-256, Keccak-256 (utilisé par Ethereum).}

\vspace{0.4cm}

\textbf{NIST} (National Institute of Standards and Technology)\\
\hspace*{1em} L'agence américaine qui définit les \textbf{standards cryptographiques} utilisés
mondialement. En 2024, le NIST a officiellement standardisé Dilithium (FIPS 204) et Falcon
(FIPS 206) comme signatures post-quantiques.

\vspace{0.4cm}

\textbf{FIPS 204 / FIPS 206}\\
\hspace*{1em} Les numéros de standards officiels du NIST. \textbf{FIPS 204} = Dilithium (ML-DSA).
\textbf{FIPS 206} = Falcon (FN-DSA). Ces standards garantissent que les implémentations sont
correctes et interopérables.

\vspace{0.4cm}

\textbf{Audit de sécurité}\\
\hspace*{1em} Une \textbf{vérification indépendante} du code par des experts externes qui cherchent
des failles de sécurité. Dans la blockchain, un audit est obligatoire avant de déployer un contrat
qui gère de l'argent réel. Un audit peut coûter de 50 000\$ à 500 000\$.

\vspace{0.4cm}

\textbf{Bug Bounty} (Prime aux bugs)\\
\hspace*{1em} Un programme public où l'équipe offre des \textbf{récompenses financières} à
quiconque trouve et signale une faille de sécurité. C'est une façon de mobiliser la communauté
pour tester le code en échange d'argent.

\vspace{0.4cm}

\textbf{zkSNARK / Preuve ZK} (Zero-Knowledge Proof)\\
\hspace*{1em} Une technique cryptographique qui permet de \textbf{prouver qu'on connaît quelque
chose sans le révéler}. Par exemple : prouver qu'une signature est valide sans exécuter la
vérification on-chain (trop chère). Très complexe à implémenter.

\end{tcolorbox}

\vspace{0.4cm}

%──────────────────────────────────────────────────────────────────────
\subsection{Termes de développement logiciel}
%──────────────────────────────────────────────────────────────────────

\begin{tcolorbox}[colback=gray!4, colframe=gray!50, breakable]

\textbf{SDK} (Software Development Kit --- Kit de Développement)\\
\hspace*{1em} Un ensemble d'outils, de bibliothèques et de documentation qui permet à
d'autres développeurs d'\textbf{intégrer facilement} une technologie dans leurs applications
sans avoir à tout réécrire de zéro.

\vspace{0.4cm}

\textbf{DApp} (Decentralized Application --- Application Décentralisée)\\
\hspace*{1em} Une application dont le backend (la logique serveur) tourne sur la blockchain
plutôt que sur un serveur centralisé. La partie frontend (interface web) reste classique.

\vspace{0.4cm}

\textbf{Proxy Pattern / Upgradeable Contract}\\
\hspace*{1em} Une technique qui permet de \textbf{mettre à jour le code} d'un smart contract
après déploiement (normalement impossible). Le ``proxy'' est un contrat fixe qui redirige
les appels vers un ``implementation'' contract qui peut être remplacé.
Avantage : corriger des bugs. Risque : introduit une centralisation.

\vspace{0.4cm}

\textbf{KAT} (Known Answer Tests --- Tests à réponse connue)\\
\hspace*{1em} Des fichiers de test fournis par le NIST contenant des \textbf{entrées et sorties
officielles} précalculées. Si notre code Python produit exactement les mêmes sorties, on est
sûr que l'implémentation est correcte.

\vspace{0.4cm}

\textbf{RPC / RPC URL}\\
\hspace*{1em} RPC (Remote Procedure Call) : le protocole utilisé pour \textbf{parler à un
noeud Ethereum}. Une RPC URL est l'adresse du serveur qui donne accès à la blockchain.
\textit{Exemples de fournisseurs : Infura, Alchemy, QuickNode.}

\vspace{0.4cm}

\textbf{Wrapper} (envelopper)\\
\hspace*{1em} Dans le contexte SuVa, wrapping = déposer un token existant (USDT) dans un
contrat et recevoir en échange un token ``enveloppé'' (sUSDT) qui représente l'original mais
avec des fonctionnalités supplémentaires. C'est réversible : on peut toujours ``unwrapper''
pour récupérer l'original.

\end{tcolorbox}

\newpage
%══════════════════════════════════════════════════════════════════════
\section{Vision du projet SuVa}
%══════════════════════════════════════════════════════════════════════

\begin{tcolorbox}[colback=suva_blue!8, colframe=suva_blue, title=\textbf{Qu'est-ce que SuVa ?}]
\textbf{SuVa (Super Value) Protocol} est un stablecoin \textbf{quantique-résistant} pour blockchains EVM.
Il permet de :
\begin{enumerate}
  \item \textbf{Wrapper} des stablecoins existants (USDT, USDC) en versions sécurisées post-quantiques
  \item Utiliser des \textbf{signatures de Dilithium ou Falcon} (NIST FIPS 204/206) à la place d'ECDSA
  \item Protéger les fonds contre les futures attaques d'ordinateurs quantiques
\end{enumerate}

\textbf{Architecture cible :}
\begin{itemize}
  \item \textbf{Off-chain (Python)} : Génération de clés, signature des transactions
  \item \textbf{On-chain (Solidity/EVM)} : Vérification des signatures, gestion du vault
  \item \textbf{Interface utilisateur} : SDK, wallet ERC-4337
\end{itemize}
\end{tcolorbox}

\begin{center}
\begin{tikzpicture}[
  node distance=1.5cm,
  box/.style={draw, rounded corners, minimum width=3cm, minimum height=0.8cm, align=center, font=\small},
  arrow/.style={-Stealth, thick, suva_blue},
]
  \node[box, fill=done_bg, draw=done_green] (python) {Python off-chain\\(keygen, sign)};
  \node[box, fill=suva_blue!10, draw=suva_blue, right=2.5cm of python] (solidity) {Solidity on-chain\\(verify, vault)};
  \node[box, fill=phase_blue, draw=suva_blue!60, above right=0.5cm and 1cm of python] (sdk) {SDK / ERC-4337\\Wallet};
  \node[box, fill=done_bg, draw=done_green, below=1.2cm of python] (ntt) {NTT, polynômes\\(fait en Solidity)};
  \node[box, fill=inprog_bg, draw=inprog_orange, below=1.2cm of solidity] (falcon) {Falcon / ETHFALCON\\(intégration)};

  \draw[arrow] (python) -- (solidity) node[midway, above, font=\tiny] {signature hex};
  \draw[arrow] (sdk) -- (python);
  \draw[arrow] (sdk) -- (solidity);
  \draw[arrow] (ntt) -- (solidity);
  \draw[arrow] (falcon) -- (solidity);
\end{tikzpicture}
\end{center}

\newpage
%══════════════════════════════════════════════════════════════════════
\section{Tableau de bord global}
%══════════════════════════════════════════════════════════════════════

\begin{center}
\renewcommand{\arraystretch}{1.4}
\begin{tabular}{|>{\bfseries}p{5.5cm}|c|p{6cm}|}
\hline
\rowcolor{suva_blue!20}
\textbf{Composant} & \textbf{Statut} & \textbf{Remarque} \\
\hline
Primitives arithmétiques (NTT) & \statdone & Solidity + Python validés par KAT NIST \\
\hline
Dilithium2/3/5 Python (keygen, sign, verify) & \statdone & Tests KAT passent \\
\hline
Dilithium2 Solidity (verify) & \statdone & ZKNOX\_dilithium2.sol opérationnel \\
\hline
ETHDilithium (variante Keccak) Python & \statdone & Optimisé pour EVM \\
\hline
ETHDilithium Solidity (verify) & \statdone & Tests croisés Python$\leftrightarrow$Solidity \\
\hline
\rowcolor{inprog_bg}
Intégration ETHFALCON (Falcon) & \stattodo & ZKNox a le code, pas encore intégré dans SuVa \\
\hline
\rowcolor{todo_bg}
Vault Contract (wrap USDT$\to$sUSDT) & \stattodo & Pas encore commencé \\
\hline
\rowcolor{todo_bg}
ERC-4337 Account Abstraction Wallet & \stattodo & Pas encore commencé \\
\hline
\rowcolor{todo_bg}
SDK Python/JS pour développeurs & \stattodo & Pas encore commencé \\
\hline
\rowcolor{todo_bg}
SPHINCS+ Solidity (fallback) & \stattodo & Très complexe, 50--500M gas \\
\hline
\rowcolor{todo_bg}
Audit de sécurité & \stattodo & Dernière étape \\
\hline
\rowcolor{todo_bg}
Déploiement mainnet & \stattodo & Post-audit \\
\hline
\end{tabular}
\end{center}

\vspace{0.5cm}
\begin{center}
\begin{tikzpicture}
  % Barre de progression globale
  \def\totaltasks{12}
  \def\donetasks{5}
  \def\barwidth{12}

  \fill[done_bg, draw=done_green, rounded corners] (0,0) rectangle (\barwidth * \donetasks / \totaltasks, 0.6);
  \fill[gray!20, draw=gray!40, rounded corners] (\barwidth * \donetasks / \totaltasks, 0) rectangle (\barwidth, 0.6);
  \node[font=\bfseries\small] at (\barwidth/2, 0.3) {\donetasks/\totaltasks\ composants terminés (41\%)};
\end{tikzpicture}
\end{center}

\newpage
%══════════════════════════════════════════════════════════════════════
\section{Ce qui est FAIT --- Détail}
%══════════════════════════════════════════════════════════════════════

\begin{done}{Couche arithmétique --- NTT et polynômes}
\textbf{Fichiers concernés :}
\begin{itemize}
  \item \texttt{python-ref/Falcon\_dilithium\_py/polynomials/polynomials.py} --- Anneaux polynomiaux $R_q$
  \item \texttt{python-ref/Falcon\_dilithium\_py/modules/modules.py} --- Modules $M_{k \times l}$
  \item \texttt{python-ref/Falcon\_dilithium\_py/utilities/utils.py} --- decompose, HighBits, LowBits, MakeHint, UseHint, check\_norm\_bound
  \item \texttt{src/ZKNOX\_NTT.sol} --- NTT en Solidity + Yul avec twiddle factors en bytecode
  \item \texttt{src/ZKNOX\_keccak\_prng.sol} --- PRNG Keccak compatible Python$\leftrightarrow$Solidity

\end{itemize}
\textbf{Tests :} 4 tests KAT NIST passent. La NTT Python et Solidity donnent des résultats identiques.\\
\textbf{Qualité :} Production-ready pour cette couche.
\end{done}

\begin{done}{Dilithium2/3/5 Python (off-chain)}
\textbf{Fichier principal :} \texttt{Falcon\_dilithium\_py/Falcon\_dilithium/Falcon\_dilithium.py}

\textbf{Ce que fait le code :}
\begin{enumerate}
  \item \texttt{keygen()} : génère $(pk, sk)$ à partir d'un seed aléatoire (SHAKE-256 PRNG)
  \item \texttt{sign(sk, message)} : signe avec rejection sampling ($\approx 4.25$ itérations)
  \item \texttt{verify(pk, message, sig)} : vérifie la signature
\end{enumerate}

\textbf{Vecteurs de test :} fichiers \texttt{PQCsignKAT\_Dilithium2/3/5.rsp} (NIST official KAT)\\
\textbf{Tailles :} pk=1312B, sk=2560B, sig=2420B (Dilithium2)
\end{done}

\begin{done}{ETHDilithium Python (variante Ethereum)}
\textbf{Fichier :} \texttt{Falcon\_dilithium\_eth.py}

\textbf{Différences avec Dilithium standard :}
\begin{itemize}
  \item PRNG \textbf{Keccak-256} à la place de SHAKE-256 (natif EVM, 3 gas vs $\sim$50\% du gas total)
  \item La clé publique inclut $\hat{A}$ (matrice précalculée en NTT) pour économiser le gas on-chain
  \item Hash du message avec Keccak-256
\end{itemize}
\textbf{Résultat :} $\sim$6.6M gas (vs $\sim$13.5M gas pour Dilithium2 standard)
\end{done}

\begin{done}{Dilithium Solidity --- Vérificateur on-chain}
\textbf{Contrats principaux :}
\begin{itemize}
  \item \texttt{src/ZKNOX\_dilithium2.sol} --- \texttt{verify(bytes pk, bytes msg, bytes sig) returns (bool)}
  \item \texttt{src/ZKNOX\_NTT.sol} --- NTT optimisée Yul pour la vérification
  \item \texttt{src/ZKNOX\_keccak\_prng.sol} --- ExpandA, ExpandMask en Keccak
\end{itemize}
\textbf{Architecture de vérification :}
\begin{enumerate}
  \item Unpack de la signature (z, h, c\_tilde)
  \item Recompute $\hat{A}$ depuis seed ou depuis pk (ETHDilithium)
  \item Calcul $w' = Az - ct_1 \cdot 2^d$
  \item UseHint pour retrouver $r_1$
  \item Hash et comparaison avec $\tilde{c}$
\end{enumerate}
\textbf{Coût gas :} $\sim$13.5M gas (Dilithium2) / $\sim$6.6M gas (ETHDilithium2)
\end{done}

\newpage
%══════════════════════════════════════════════════════════════════════
\section{Ce qui RESTE À FAIRE --- Tâches prioritaires}
%══════════════════════════════════════════════════════════════════════

%──────────────────────────────────────────────────────────────────────
\subsection{Priorité CRITIQUE : Intégration Falcon / ETHFALCON}
%──────────────────────────────────────────────────────────────────────

\begin{critical}{Falcon --- La cible production de SuVa}
\textbf{Pourquoi c'est critique :}\\
Dilithium coûte $\sim$6--13M gas par vérification. C'est \textbf{100x trop cher} pour un usage en production. Falcon (NIST FIPS 206) coûte seulement $\sim$\textbf{350K gas}.

\textbf{Ce qui existe déjà (ZKNox) :}
\begin{itemize}
  \item \texttt{ETHFALCON} : implémentation Falcon pour EVM par ZKNox
  \item Repo GitHub : \texttt{github.com/ZKNox/ETHFALCON}
\end{itemize}

\textbf{Ce qui reste à faire pour SuVa :}
\begin{enumerate}
  \item \textbf{Étudier ETHFALCON} : comprendre les contrats Solidity de vérification
  \item \textbf{Adapter ETHFALCON} au contexte SuVa (interface vault, ERC-4337)
  \item \textbf{Python off-chain} : implémenter ou adapter le signing Falcon en Python
  \item \textbf{Tests croisés} : vecteurs Python $\to$ Solidity qui passent
  \item \textbf{Intégrer dans le vault} : le vault accepte les signatures Falcon
\end{enumerate}

\textbf{Comparaison Falcon vs Dilithium :}
\begin{center}
\begin{tabular}{lrrrr}
\toprule
\textbf{Schéma} & \textbf{Gas verify} & \textbf{Sig taille} & \textbf{PK taille} & \textbf{Sécurité}\\
\midrule
Dilithium2      & $\sim$13.5M & 2420 B & 1312 B & 128 bits classique \\
ETHDilithium2   & $\sim$6.6M  & 2420 B & >1312 B & 128 bits classique \\
\rowcolor{done_bg}
Falcon-512      & $\sim$350K  &  666 B &  897 B & 128 bits classique \\
SPHINCS+        & 50--500M    & 8--50 KB & 32 B & 128 bits (hash)\\
\bottomrule
\end{tabular}
\end{center}
\end{critical}

%──────────────────────────────────────────────────────────────────────
\subsection{Priorité HAUTE : Vault Contract}
%──────────────────────────────────────────────────────────────────────

\begin{todo}{Vault Contract --- Wrapping de stablecoins}
\textbf{Rôle :} Contrat Solidity qui verrouille des USDT/USDC et émet des sUSDT/sUSDC.

\textbf{Fonctionnalités requises :}
\begin{enumerate}
  \item \texttt{deposit(amount, pk\_pq)} : dépose USDT, stocke la clé publique PQ
  \item \texttt{withdraw(amount, sig\_pq, msg)} : retire USDT après vérification de la signature PQ
  \item \texttt{transfer(to, amount, sig\_pq, msg)} : transfert sUSDT après vérification PQ
  \item Gestion des nonces pour éviter les replay attacks
  \item Interface avec les contrats de vérification Dilithium/Falcon
\end{enumerate}

\textbf{Questions ouvertes :}
\begin{itemize}
  \item Quel standard ERC pour sUSDT ? ERC-20 avec custom transfer ? ERC-4337 ?
  \item Comment gérer la mise à jour des clés PQ (rotation de clés) ?
  \item Comment gérer le cas où l'utilisateur perd sa clé PQ ?
\end{itemize}

\textbf{Estimation effort :} 2--4 semaines de développement + tests
\end{todo}

%──────────────────────────────────────────────────────────────────────
\subsection{Priorité HAUTE : ERC-4337 Account Abstraction}
%──────────────────────────────────────────────────────────────────────

\begin{todo}{ERC-4337 --- Wallet quantique-résistant}
\textbf{Rôle :} Remplacer ECDSA par Dilithium/Falcon dans les smart contract wallets.

\textbf{Fonctionnement ERC-4337 :}
\begin{enumerate}
  \item L'utilisateur soumet une \texttt{UserOperation} (au lieu d'une transaction ECDSA)
  \item Un \texttt{Bundler} la regroupe avec d'autres et l'envoie au \texttt{EntryPoint} contract
  \item L'\texttt{EntryPoint} appelle \texttt{validateUserOp()} sur le wallet de l'utilisateur
  \item Le wallet appelle le vérificateur Dilithium/Falcon pour valider la signature PQ
\end{enumerate}

\textbf{Fichiers à créer :}
\begin{itemize}
  \item \texttt{src/SuVaWallet.sol} : smart contract wallet ERC-4337 avec PQ validation
  \item \texttt{src/SuVaPaymaster.sol} (optionnel) : pour payer le gas en sUSDT
  \item \texttt{python-ref/wallet\_sdk.py} : SDK pour créer des UserOperations
\end{itemize}

\textbf{Défi technique :} Le gas de validation (Dilithium = 6.6M gas) est 100x le coût ECDSA (ecrecover = 3000 gas). C'est pour cela que Falcon ($\sim$350K gas) est essentiel.
\end{todo}

%──────────────────────────────────────────────────────────────────────
\subsection{Priorité MOYENNE : SDK}
%──────────────────────────────────────────────────────────────────────

\begin{todo}{SDK --- Interface pour les développeurs}
\textbf{Composants :}
\begin{enumerate}
  \item \textbf{Python SDK} : wrapper de haut niveau pour keygen, sign, broadcast transaction
  \item \textbf{JavaScript/TypeScript SDK} : pour les DApps web
  \item \textbf{Documentation API} : guide d'intégration pour d'autres projets
\end{enumerate}

\textbf{Fonctions clés du SDK :}
\begin{itemize}
  \item \texttt{generate\_keypair(scheme='falcon512')} $\to$ (pk, sk)
  \item \texttt{sign\_transaction(sk, tx\_data)} $\to$ signature
  \item \texttt{create\_user\_operation(pk, sk, tx\_data)} $\to$ UserOp (ERC-4337)
  \item \texttt{broadcast(user\_op, rpc\_url)} $\to$ tx\_hash
\end{itemize}
\end{todo}

%──────────────────────────────────────────────────────────────────────
\subsection{Priorité BASSE : SPHINCS+ (fallback)}
%──────────────────────────────────────────────────────────────────────

\begin{todo}{SPHINCS+ --- Le ``bunker nucléaire''}
\textbf{Rôle :} Schéma de signature basé uniquement sur des fonctions de hachage (pas de réseaux de treillis).
Si une attaque contre Falcon/Dilithium était découverte, SPHINCS+ reste sûr.

\textbf{Problème majeur :} 50 à 500 millions de gas par vérification sur EVM. Pratiquement inutilisable.

\textbf{Options possibles :}
\begin{itemize}
  \item Vérification off-chain avec preuve ZK (zkSNARK) soumise on-chain --- très complexe
  \item Layer 2 spécialisé avec SPHINCS+ natif --- architecture différente
  \item Utilisation uniquement pour les opérations non-urgentes (governance, key rotation)
\end{itemize}

\textbf{Décision :} Repousser jusqu'après Falcon. ZKNox n'a pas d'implémentation EVM de SPHINCS+.
\end{todo}

\newpage
%══════════════════════════════════════════════════════════════════════
\section{Feuille de route --- Planning estimé}
%══════════════════════════════════════════════════════════════════════

\begin{center}
\begin{tikzpicture}[
  x=1.4cm, y=1.2cm,
  phase/.style={draw, rounded corners, minimum height=0.7cm, align=center, font=\small\bfseries},
  done_style/.style={phase, fill=done_bg, draw=done_green},
  inprog_style/.style={phase, fill=inprog_bg, draw=inprog_orange},
  todo_style/.style={phase, fill=todo_bg, draw=todo_red!60},
  label_style/.style={font=\footnotesize\bfseries, align=center},
]
  % Axe du temps
  \draw[-Stealth, thick, suva_blue] (0, 0) -- (9.5, 0) node[right] {\small Temps};

  % Phases
  \node[done_style, minimum width=3.5cm] at (1.7, 1.5) {Phase 1\\Primitives + Dilithium};
  \node[inprog_style, minimum width=2cm] at (4.5, 1.5) {Phase 2\\Falcon};
  \node[todo_style, minimum width=2.5cm] at (6.7, 1.5) {Phase 3\\Vault + 4337};
  \node[todo_style, minimum width=2cm] at (8.8, 1.5) {Phase 4\\SDK + Audit};

  % Marqueurs
  \foreach \x/\label in {0/Déc 2024, 3.5/Mars 2026, 5.5/Mai 2026, 8/Juil 2026, 9.5/Oct 2026} {
    \draw (\x, -0.15) -- (\x, 0.15);
    \node[below, font=\tiny, rotate=20] at (\x, -0.2) {\label};
  }

  % Statuts
  \node[done_style, font=\tiny, minimum width=1.5cm] at (1.7, 0.5) {TERMINÉ};
  \node[inprog_style, font=\tiny, minimum width=1cm] at (4.5, 0.5) {ACTUEL};
  \node[todo_style, font=\tiny, minimum width=1.5cm] at (6.7, 0.5) {À FAIRE};
  \node[todo_style, font=\tiny, minimum width=1cm] at (8.8, 0.5) {À FAIRE};
\end{tikzpicture}
\end{center}

\vspace{0.5cm}

\begin{center}
\renewcommand{\arraystretch}{1.5}
\begin{tabular}{|c|p{4cm}|p{7cm}|c|}
\hline
\rowcolor{suva_blue!20}
\textbf{Phase} & \textbf{Nom} & \textbf{Livrables} & \textbf{Durée} \\
\hline
\rowcolor{done_bg}
1 & \textbf{Fondations} (FAIT) &
NTT Solidity, Dilithium Python, Vérificateur Solidity, Tests KAT &
$\sim$3 mois \\
\hline
\rowcolor{inprog_bg}
2 & \textbf{Falcon} (PRIORITÉ) &
Étude ETHFALCON, Python off-chain, Tests croisés, Intégration &
8--10 sem. \\
\hline
\rowcolor{todo_bg}
3a & \textbf{Vault Contract} &
Dépôt USDT$\to$sUSDT, Retrait, Transfer avec sig PQ, Nonces &
4--6 sem. \\
\hline
\rowcolor{todo_bg}
3b & \textbf{ERC-4337 Wallet} &
Smart wallet, validateUserOp, Bundler, Paymaster (optionnel) &
4--6 sem. \\
\hline
\rowcolor{todo_bg}
4a & \textbf{SDK} &
Python SDK, JS/TS SDK, Documentation, Exemples d'intégration &
3--4 sem. \\
\hline
\rowcolor{todo_bg}
4b & \textbf{Audit + Mainnet} &
Audit de sécurité externe, Bug bounty, Déploiement Ethereum mainnet &
$\sim$2 mois \\
\hline
\end{tabular}
\end{center}

\newpage
%══════════════════════════════════════════════════════════════════════
\section{Tâches immédiates pour Mohamed Lamine MBENGUE}
%══════════════════════════════════════════════════════════════════════

\begin{tcolorbox}[colback=suva_blue!5, colframe=suva_blue, title=\textbf{Prochaines actions concrètes (semaines 1--2)}]
\begin{enumerate}[leftmargin=*]
  \item[\textbf{J1}] \textbf{Installer l'environnement et faire passer les tests}
    \begin{itemize}
      \item Créer le venv Python, installer les dépendances
      \item Lancer \texttt{pytest tests/ -v} $\to$ 4 tests VERTS
      \item Installer Foundry et faire \texttt{forge build}
    \end{itemize}

  \item[\textbf{J2--J3}] \textbf{Maîtriser le code Dilithium}
    \begin{itemize}
      \item Lire et annoter \texttt{Falcon\_dilithium.py} ligne par ligne
      \item Faire tous les exercices du Chapitre 8 (ce document)
      \item Comprendre \texttt{ZKNOX\_dilithium2.sol} et son architecture
    \end{itemize}

  \item[\textbf{J4--J5}] \textbf{Explorer ETHFALCON (ZKNox)}
    \begin{itemize}
      \item Cloner \texttt{github.com/ZKNox/ETHFALCON}
      \item Lire le README et les contrats Solidity
      \item Identifier ce qui est réutilisable pour SuVa
    \end{itemize}

  \item[\textbf{Semaine 2}] \textbf{Préparer les questions pour le client}
    \begin{itemize}
      \item Architecture Vault : ERC-20 standard ou ERC-4337 d'emblée ?
      \item Priorité : Falcon seul, ou Dilithium + Falcon en parallèle ?
      \item Cible réseau : Ethereum mainnet, ou d'abord un L2 (Polygon, Arbitrum) ?
      \item Gestion de la perte de clé : quel mécanisme de récupération ?
      \item Budget gas acceptable en production : < 500K gas ? < 1M gas ?
    \end{itemize}
\end{enumerate}
\end{tcolorbox}

%══════════════════════════════════════════════════════════════════════
\section{Questions à poser au client}
%══════════════════════════════════════════════════════════════════════

\begin{tcolorbox}[colback=warn_bg, colframe=orange!80!black,
  title=\textbf{Questions techniques prioritaires}, breakable]

\textbf{Sur l'architecture :}
\begin{enumerate}
  \item L'approche retenue est-elle Falcon uniquement, ou Dilithium (dev/test) + Falcon (prod) ?
  \item Le vault doit-il être upgradeable (proxy pattern) ou immutable ?
  \item Quelle est la cible réseau principale ? Ethereum mainnet ? Un L2 ?

\end{enumerate}

\textbf{Sur Falcon / ETHFALCON :}
\begin{enumerate}[resume]
  \item Avez-vous déjà contacté ZKNox pour l'intégration d'ETHFALCON ?
  \item Y a-t-il des modifications prévues à ETHFALCON ou on l'utilise tel quel ?
  \item Comment gérer la génération de clé Falcon (qui requiert un sampling gaussien discret) ?
\end{enumerate}

\textbf{Sur l'usage et l'UX :}
\begin{enumerate}[resume]
  \item Quel est le scénario d'usage principal : wallets hardware, mobile, ou serveurs ?
  \item Comment l'utilisateur stocke-t-il sa clé secrète Falcon/Dilithium ?
  \item Quelle est la stratégie si un utilisateur perd sa clé privée PQ ?
\end{enumerate}

\textbf{Sur le budget et les délais :}
\begin{enumerate}[resume]
  \item Quel est le budget gas maximum acceptable par transaction ?
  \item Y a-t-il une date cible pour un déploiement testnet ? Mainnet ?
  \item Un audit de sécurité externe est-il prévu, et par qui ?
\end{enumerate}
\end{tcolorbox}

%══════════════════════════════════════════════════════════════════════
\section{Risques et points de vigilance}
%══════════════════════════════════════════════════════════════════════

\begin{center}
\renewcommand{\arraystretch}{1.4}
\begin{tabular}{|p{4.5cm}|c|c|p{5cm}|}
\hline
\rowcolor{suva_blue!20}
\textbf{Risque} & \textbf{Probabilité} & \textbf{Impact} & \textbf{Mitigation} \\
\hline
ETHFALCON instable ou non audité & Moyen & Critique & Contribuer à l'audit, tests exhaustifs \\
\hline
Gas Falcon > 350K en pratique & Faible & Haut & Profiler avec forge --gas-report \\
\hline
Perte de clé PQ par utilisateur & Haut & Critique & Mécanisme de récupération sociale ou multi-sig \\
\hline
Évolution du standard FIPS 204/206 & Faible & Moyen & Suivre les errata NIST \\
\hline
Attack on Falcon sampling (Gaussian) & Très faible & Critique & Utiliser l'implémentation de référence \\
\hline
Pas de support EIP pour les ops PQ & Moyen & Moyen & ERC-4337 contourne ce problème \\
\hline
\end{tabular}
\end{center}

\vspace{1cm}
\begin{center}
  {\color{suva_blue}\rule{0.5\linewidth}{1pt}}\\[0.3cm]
  {\small\color{gray} Document de suivi de projet créé par \textbf{Mohamed Lamine MBENGUE} ---
  Mars 2026 ---
  \href{https://portfolio-alamine.web.app}{\texttt{portfolio-alamine.web.app}}}
\end{center}

\end{document}
